\section{Introduction}\label{sec:introduction}

Lieb--Thirring inequalities bound moments of the negative eigenvalues of a Schr\"odinger operator by an integral of its potential. Their sharp constants measure how much binding is possible for a prescribed potential cost. When the wave function has several components, the potential becomes a matrix, and its eigenspaces may change with position. The question addressed here is whether this additional freedom increases the optimal constant in one dimension.

Let $m\ge1$, let $1/2<\gamma<3/2$, and put $p=\gamma+1/2$.
For a measurable Hermitian positive semidefinite matrix function
$W:\R\to\C^{m\times m}$ satisfying
\begin{equation}\label{eq:potential-class}
  \int_\R\tr(W(x)^p)\,dx<\infty,
\end{equation}
consider the quadratic form
\begin{equation}\label{eq:schrodinger-form}
 h_W[\psi]=\int_\R\|\psi'(x)\|^2\,dx
       -\int_\R\langle\psi(x),W(x)\psi(x)\rangle\,dx,
 \qquad \psi\in H^1(\R;\C^m).
\end{equation}
This form is closed and bounded below. Its self-adjoint operator is denoted by
$H_W=-d^2/dx^2-W$. Its negative spectrum consists of eigenvalues, with finite
multiplicities and possible accumulation only at zero; these facts are proved
under exactly \eqref{eq:potential-class} in Section~\ref{sec:spectral}.
We write
\[
 \Tr(H_W)_-^\gamma=\sum_{\lambda_j(H_W)<0}|\lambda_j(H_W)|^\gamma,
\]
initially allowing the sum to be infinite. Matrix powers are defined by spectral
functional calculus, and $\tr$ is the ordinary, unnormalized matrix trace.

\begin{theorem}[Sharp matrix inequality]\label{thm:main}
For every $1/2<\gamma<3/2$, every finite $m\ge1$, and every $W$ satisfying
\eqref{eq:potential-class},
\begin{equation}\label{eq:main-bound}
 \Tr(H_W)_-^\gamma\le C_\gamma\int_\R\tr(W(x)^{\gamma+1/2})\,dx,
 \qquad
 C_\gamma=
 \left(\frac{\gamma-1/2}{\gamma+1/2}\right)^{\gamma-1/2}
 \frac{\Gamma(\gamma+1)}{\sqrt\pi\,\Gamma(\gamma+3/2)}.
\end{equation}
The constant is optimal in every matrix dimension. With
$r=(\gamma-1/2)^{-1}$, it is attained by
\begin{equation}\label{eq:extremal-potential}
 W(x)=\diag\big((r+1)\sech^2(rx),0,\ldots,0\big).
\end{equation}
\end{theorem}

The matrices $W(x)$ may be noncommuting and may have arbitrary rank. In
particular, the conclusion concerns changing internal channels as well as
independent scalar channels.

\subsection{Historical context and relationship to prior work}
Lieb and Thirring introduced their eigenvalue inequalities in the study
of fermionic kinetic energy and the stability of matter
\cite{LiebThirring1975,LiebThirring1976}.
The one-state variational problem goes back to Keller \cite{Keller1961}.
In one dimension its optimal constant exceeds the semiclassical constant
for $1/2<\gamma<3/2$. The scalar conjecture in this interval asks whether
summing all negative eigenvalues changes that one-state constant.
The companion paper \cite[Theorem~1.1]{ScalarCompanion} proves that it does
not. The present paper establishes the same constant for arbitrary finite
matrix size, allowing the eigenspaces of the potential to vary with position.

Sharp matrix inequalities at the two endpoint exponents have an
established history. Weidl proved finiteness at the scalar lower endpoint
\cite{Weidl1996}, and Hundertmark, Lieb and Thomas obtained the sharp
$\gamma=1/2$ inequality with constant $1/2$
\cite[Theorem~1]{HLT1998}; Hundertmark, Laptev and Weidl proved the
operator-valued version \cite[Theorem~3.1]{HLW2000}.
At $\gamma=3/2$, Laptev and Weidl established the sharp matrix estimate by
trace identities and extended it to operator-valued potentials
\cite[Theorems~2.1 and~2.2]{LaptevWeidl2000}.
Aizenman--Lieb integration in the spectral parameter gives the larger
exponents \cite{AizenmanLieb1978}. Laptev and Weidl's dimension-lifting
argument then yields the semiclassical constant in every spatial dimension
for $\gamma\ge3/2$ \cite[Theorem~3.1]{LaptevWeidl2000}.
Benguria and Loss gave a proof of the matrix upper-endpoint estimate by
Riccati factorization and commutation \cite{BenguriaLoss2000}.

In the open interval, Hundertmark, Laptev and Weidl obtained the
operator-valued upper bound with twice the semiclassical constant
\cite[Theorem~4.1]{HLW2000}.
Dolbeault, Laptev and Loss developed the method of orthonormal
vector-valued functions at exponent one
\cite[Theorems~1 and~3]{DLL2008}.
More recently, Read and Schulz proved the sharp value
$4/(3\sqrt3\pi)$ for $\gamma=1$, including operator-valued potentials
\cite[Proposition~4 and Equation~(2.1)]{ReadSchulz2026}.
Thus the exponent-one case of our theorem is already contained in their
result. Our additional range is $1/2<\gamma<3/2$, $\gamma\ne1$.

The proof adapts the finite-interval action and continuation framework
of \cite[Theorem~2.1, Proposition~5.1 and Section~6]{ScalarCompanion}.
The companion's comparison applies when the pointwise density has rank
at most one \cite[Proposition~4.1]{ScalarCompanion}.
The new analytic step removes that restriction, making the action method
applicable to vector-valued wave functions.



\subsection{The matrix comparison and the proof strategy}
The proof uses the finite-interval action method of the scalar companion
\cite{ScalarCompanion}. To explain the extension, choose finitely many
orthonormal real vector-valued trial functions $u_1,\ldots,u_n$ and place them in
the rows of a matrix function $U(x)$. A lower triangular coefficient matrix
$C$ produces $A(x)=CU(x)$, so that the $i$th row of $A$ remains in the span
of $u_1,\ldots,u_i$. For trial functions ordered by their negative energies,
this preserves the spectral inequalities needed to bound the action from above.
The action subtracts a power of the matrix density
$\rho(x)=A(x)A(x)^{\mathsf T}$, and a complementary lower bound recovers
the sum of the chosen eigenvalue moments. A fixed realification of the
potential will subsequently give the complex Hermitian case.

The lower bound is obtained by connecting two prescribed diagonal matrices
$K$ and $-K$, where the diagonal entries of $K$ are the positive square roots
of the absolute values of the negative trial energies. A regularized matrix field $M$ drives an equation with the
penalized constraint $M(B)=\rho$. In the scalar-potential argument,
$A$ has one column and this constraint has rank at most one. For matrix
potentials, $\rho$ can have any rank. The essential new estimate controls
$\tr(\rho^q)$ in this setting, where $q$ is conjugate to $p$.

We prove that estimate as an independent result in Section~\ref{sec:trace}.
For Hermitian $B,V$ such that $(sI-B)^2+V$ is uniformly positive for real $s$,
we form a weighted integral $N$ of resolvent differences. When $N\ge0$,
the result is $\tr(NV)\ge\tr(N^q)$. The hypothesis permits indefinite $V$,
which is necessary when $V$ is later obtained from $K^2-B^2$.
Half-line solutions of a constant-coefficient differential equation give
two boundary matrices. Their sum identifies the integrated resolvent,
and their difference accounts for the noncommutative error. A logarithmic
determinant identity and a nonnegative matrix square then prove the trace
comparison with the required constant.

Section~\ref{sec:action-framework} states the action inequality and isolates
its relation to spectral estimates. Sections~\ref{sec:field}--\ref{sec:action-limit}
construct the field, connect the endpoints, and remove the penalty. The
continuation argument works modulo independent signs of the rows of $C$.
After smoothing and perturbation, parity of the boundary of the quotient
zero set forces an endpoint solution. Finally,
Section~\ref{sec:spectral} applies the action inequality to smooth trial
spaces and passes to general potentials and the full spectral moment.
All of these steps are proved here; the companion is cited for the method
and its scalar specialization.

\subsection{Conventions}
For the rest of the paper we use
\begin{equation}\label{eq:exponents}
 \sigma=\gamma-\tfrac12\in(0,1),\qquad
 p=1+\sigma,\qquad q=1+\sigma^{-1},\qquad
 \beta=\sigma-\tfrac12.
\end{equation}
Thus $p^{-1}+q^{-1}=1$. The identity matrix has the size specified by its
context. On matrices we use the Hilbert--Schmidt norm
$\|A\|_{\HS}^2=\tr(A^*A)$ and the operator norm $\|A\|_{\op}$.
The space $\Sym_n$ of real symmetric matrices carries the inner product
$\tr(BH)$. Complex inner products are linear in the second variable.
